
\documentclass[aps,preprint]{revtex4}
%%%%%%%%%%%%%%%%%%%%%%%%%%%%%%%%%%%%%%%%%%%%%%%%%%%%%%%%%%%%%%%%%%%%%%%%%%%%%%%%%%%%%%%%%%%%%%%%%%%%%%%%%%%%%%%%%%%%%%%%%%%%%%%%%%%%%%%%%%%%%%%%%%%%%%%%%%%%%%%%%%%%%%%%%%%%%%%%%%%%%%%%%%%%%%%%%%%%%%%%%%%%%%%%%%%%%%%%%%%%%%%%%%%%%%%%%%%%%%%%%%%%%%%%%%%%
\usepackage{amsfonts}
\usepackage{amsmath}
\usepackage{amssymb}
\usepackage{graphicx}

\setcounter{MaxMatrixCols}{10}
%TCIDATA{OutputFilter=LATEX.DLL}
%TCIDATA{Version=5.50.0.2960}
%TCIDATA{<META NAME="SaveForMode" CONTENT="1">}
%TCIDATA{BibliographyScheme=Manual}
%TCIDATA{Created=Friday, April 06, 2012 11:15:25}
%TCIDATA{LastRevised=Thursday, June 26, 2025 14:20:45}
%TCIDATA{<META NAME="GraphicsSave" CONTENT="32">}
%TCIDATA{<META NAME="DocumentShell" CONTENT="Articles\SW\REVTeX 4">}
%TCIDATA{Language=American English}
%TCIDATA{CSTFile=revtex4.cst}

\newtheorem{theorem}{Theorem}
\newtheorem{acknowledgement}[theorem]{Acknowledgement}
\newtheorem{algorithm}[theorem]{Algorithm}
\newtheorem{axiom}[theorem]{Axiom}
\newtheorem{claim}[theorem]{Claim}
\newtheorem{conclusion}[theorem]{Conclusion}
\newtheorem{condition}[theorem]{Condition}
\newtheorem{conjecture}[theorem]{Conjecture}
\newtheorem{corollary}[theorem]{Corollary}
\newtheorem{criterion}[theorem]{Criterion}
\newtheorem{definition}[theorem]{Definition}
\newtheorem{example}[theorem]{Example}
\newtheorem{exercise}[theorem]{Exercise}
\newtheorem{lemma}[theorem]{Lemma}
\newtheorem{notation}[theorem]{Notation}
\newtheorem{problem}[theorem]{Problem}
\newtheorem{proposition}[theorem]{Proposition}
\newtheorem{remark}[theorem]{Remark}
\newtheorem{solution}[theorem]{Solution}
\newtheorem{summary}[theorem]{Summary}

\newenvironment{proof}[1][Proo

\title{Summary of Discussion: Nuclear Operators, Tensor Products, and Operator Topologies}
\author{DeepSeek Chat}
\date{\today}

\newcommand{\HH}{\mathcal{H}}
\newcommand{\BB}{\mathcal{B}}
\newcommand{\KK}{\mathcal{K}}
\newcommand{\LL}{\mathcal{L}}
\newcommand{\ot}{\otimes}

\begin{document}

\maketitle

\section{Key Definitions and Initial Setup}
Let \(\HH\) be an infinite-dimensional complex Hilbert space. We consider:
\begin{itemize}
    \item \textbf{Nuclear (Trace-Class) Operators}: \(\LL^1(\HH)\), with norm \(\|T\|_1 = \sum_i \sigma_i\) (singular values).
    \item \textbf{Hilbert-Schmidt Operators}: \(\LL^2(\HH)\), with norm \(\|T\|_2 = \left(\sum_i \sigma_i^2\right)^{1/2}\).
    \item \textbf{Compact Operators}: \(\KK(\HH)\), norm-closure of finite-rank operators.
    \item \textbf{Bounded Operators}: \(\BB(\HH)\), all continuous linear maps \(\HH \to \HH\).
\end{itemize}

\section{Tensor Products and Operator Spaces}
The algebraic tensor product \(\HH \ot \HH^*\) (mixed tensors) corresponds to finite-rank operators via:
\[
x \ot y^* \mapsto \left( z \mapsto \langle y, z \rangle x \right).
\]
Different completions yield distinct operator spaces:
\begin{table}[h]
    \centering
    \begin{tabular}{|l|l|l|}
        \hline
        \textbf{Norm on \(\HH \ot \HH^*\)} & \textbf{Completion} & \textbf{Resulting Space} \\ \hline
        Operator norm (\(\|\cdot\|_{\text{op}}\)) & \(\HH \ot_\epsilon \HH^*\) & \(\KK(\HH)\) (compact operators) \\ \hline
        Trace norm (\(\|\cdot\|_1\)) & \(\HH \ot_\pi \HH^*\) & \(\LL^1(\HH)\) (trace-class) \\ \hline
        Hilbert-Schmidt norm (\(\|\cdot\|_2\)) & \(\HH \ot_2 \HH^*\) & \(\LL^2(\HH)\) (Hilbert-Schmidt) \\ \hline
        Strong Operator Topology (SOT) & SOT-closure & \(\BB(\HH)\) (all bounded operators) \\ \hline
    \end{tabular}
    \caption{Completions of \(\HH \ot \HH^*\) under different topologies.}
\end{table}

\section{Strong Operator Topology (SOT) and Density}
\begin{theorem}[SOT-Density of Finite-Rank Operators]
The space \(\HH \ot \HH^*\) (finite-rank operators) is \textbf{dense} in \(\BB(\HH)\) under the SOT. Thus:
\[
\overline{\HH \ot \HH^*}^{\text{SOT}} = \BB(\HH).
\]
\end{theorem}
\begin{proof}
For any \(T \in \BB(\HH)\) and finite set \(\{x_1, \dots, x_n\} \subset \HH\), there exists a finite-rank operator \(T_0\) such that:
\[
\|(T - T_0)x_i\| < \epsilon \quad \forall i.
\]
This follows by projecting \(T\) onto \(\operatorname{span}\{x_1, \dots, x_n\}\).
\end{proof}

\section{Isomorphism Questions}
\subsection{Algebraic Isomorphism}
At the algebraic level:
\[
\HH \ot \HH^* \cong \text{Finite-rank operators} \subset \BB(\HH).
\]

\subsection{Topological Isomorphisms}
\begin{itemize}
    \item \textbf{Uniform (Operator Norm)}: \(\overline{\HH \ot \HH^*} = \KK(\HH) \subsetneq \BB(\HH)\).
    \item \textbf{SOT}: \(\overline{\HH \ot \HH^*}^{\text{SOT}} = \BB(\HH)\), but this is \textbf{not} a Banach space isomorphism (SOT is not normable).
    \item \textbf{Trace/Hilbert-Schmidt Norms}: Completions yield proper subspaces (\(\LL^1(\HH), \LL^2(\HH)\)).
\end{itemize}

\section{Conclusion}
\begin{itemize}
    \item The SOT-completion of \(\HH \ot \HH^*\) covers \textbf{all} bounded operators, but only as a \textbf{topological vector space} (no norm).
    \item Normed completions (\(\|\cdot\|_{\text{op}}, \|\cdot\|_1, \|\cdot\|_2\)) restrict to smaller subspaces.
    \item The isomorphism \(\HH \ot \HH^* \to \BB(\HH)\) is \textbf{algebraically dense} and \textbf{SOT-continuous}, but lacks a normed structure.
\end{itemize}

\end{document}